%% AQPS -- Electric Power Systems Research submission (elsarticle).
%% Every number in the text is a macro from numbers.tex (generated by postprocess/make_numbers.py in the
%% simulation repository); every table is generated by postprocess/make_tables.py.
%% All files sit in one folder (Editorial Manager does not accept subfolders).
%% Layout switch: \reviewlayouttrue gives the EPSR single-column, double-spaced, 12 pt review layout
%% (elsarticle preprint + epsr.sty); \reviewlayoutfalse gives the two-column 3p layout.
\newif\ifreviewlayout
\reviewlayoutfalse
\ifreviewlayout
  \documentclass[preprint,12pt]{elsarticle}
  \usepackage{epsr}
  \usepackage{etoolbox}% tables were sized for the 3p text width; use a smaller font in the 150 mm layout
  \AtBeginEnvironment{table}{\let\footnotesize\scriptsize}
  \AtBeginEnvironment{table*}{\let\footnotesize\scriptsize\setlength{\tabcolsep}{3pt}}
\else
  \documentclass[final,3p,times,twocolumn]{elsarticle}
\fi
\usepackage{amsmath,amssymb,amsthm}
\newtheorem{proposition}{Proposition}
\usepackage{booktabs}
\usepackage{graphicx}
\usepackage{algorithm}
\usepackage{algpseudocode}
\usepackage{xcolor}
\usepackage[hidelinks]{hyperref}

\input{numbers.tex}

% float placement: keep figures and tables near their first mention
\renewcommand{\topfraction}{0.9}
\renewcommand{\bottomfraction}{0.8}
\renewcommand{\textfraction}{0.07}
\renewcommand{\floatpagefraction}{0.8}
\renewcommand{\dbltopfraction}{0.85}
\renewcommand{\dblfloatpagefraction}{0.8}
\setcounter{topnumber}{3}
\setcounter{dbltopnumber}{2}

\newcommand{\Ses}{E^{\mathrm{ses}}}
\newcommand{\Mis}{E^{\mathrm{mis}}}
\newcommand{\Psite}{P_{\mathrm{site}}}
\newcommand{\Cset}{\mathcal{C}_k}
\newcommand{\Sset}{\mathcal{S}_k}
\newcommand{\xpri}{x^{\mathrm{pri}}}

\journal{Electric Power Systems Research}

\begin{document}

\begin{frontmatter}

\title{Schedule-Aware Real-time Depot Charging for Electric Fleets}

\author[inst1]{Mohammad Kamar Uddin\corref{cor1}}
\ead{s3955218@student.rmit.edu.au}
\cortext[cor1]{Corresponding author.}
\author[inst1]{Ali Moradi Amani}
\ead{ali.moradiamani@rmit.edu.au}
\author[inst2]{Hui Song}
\ead{hui.song1@newcastle.edu.au}

\affiliation[inst1]{organization={School of Engineering, RMIT University},
            city={Melbourne}, state={VIC}, postcode={3000}, country={Australia}}
\affiliation[inst2]{organization={School of Information and Physical Sciences, University of Newcastle},
            city={Callaghan}, state={NSW}, postcode={2308}, country={Australia}}

\begin{abstract}
Commercial fleet depots charge more vehicles than they have chargers, under a transformer limit, while
vehicles must leave with enough energy for rostered trips that change during the day. This paper proposes
Adaptive Queuing-based Predictive Scheduling (AQPS), a three-layer controller. Layer~0 derives each
vehicle's mission energy from the roster, ranks vehicles by the Apparent Tardiness Cost index and admits
the largest priority set for which the Layer-2 quadratic program can meet every mission as a hard
constraint. Layer~1 allocates chargers: it preempts non-priority vehicles only when a waiting priority
vehicle's slack is exhausted, releases a priority vehicle only once its mission energy is delivered, and
otherwise escalates to the operator. Layer~2 solves a convex receding-horizon quadratic program only on
events and feeds predicted shortfalls back to Layer~0; an operator impact preview reports the cost of
alternatives. In simulations of 45- and 90-vehicle depots with real AEMO rooftop-PV forecasts and actuals,
AQPS met all missions in a peak-stress scenario, against \RMissionLLFSThreeFortyFive\,\% and
\RMissionLLFSThreeNinety\,\% for least-laxity-first, at an equal or lower cost per kWh; the price is lower
non-priority energy (\RNonPriRTETSThreeFortyFive\,\% vs \RNonPriLLFSThreeFortyFive\,\%). Event triggering
cut optimiser solves by \RTotRedMin--\RTotRedMax\,\% with no change in mission sufficiency.
\end{abstract}

\begin{keyword}
EV fleet charging \sep depot charging \sep priority scheduling \sep model predictive control
\sep event-triggered control \sep human-in-the-loop
\end{keyword}

\end{frontmatter}

% =============================================================================
\section{Introduction}
\label{sec:intro}

Electric light commercial vehicles are a fast-growing segment of transport electrification, driven by
net-zero mandates for fleets \cite{iea2025global,alhanahi2024charging}. Fleet vehicles run on rosters: a van
that leaves under-charged causes a service failure, not just an energy shortfall. Depot charging is
constrained on three sides: chargers (EVSEs) are fewer than vehicles, the transformer caps total power and
is expensive to upgrade \cite{motlagh2026review,keshavarzian2024optimal}, and the roster changes during the
day through breakdowns, absences and route swaps \cite{alhanahi2024charging,jack2025estimating,zhao2023toward}.
Time-of-use (ToU) tariffs and on-site PV and storage add a cost dimension \cite{corsetti2024commercial,chen2025optimal}.

Bus depots with fixed timetables are scheduled offline \cite{he2020optimal,hosseini2023optimizing,meng2024research,shi2022energy};
delivery fleets have been treated with bi-level formulations \cite{skugor2016bilevel}, with on-site PV and
storage \cite{dicorato2019integrated,vandermeer2018energy,chen2025optimal}, and with evolutionary
\cite{liu2023multiobjective}, stochastic convex \cite{qureshi2024multiobjective} and reinforcement-learning
\cite{orfanoudakis2025ev2gym,mei2026multiagent,vales2026emergent} schedulers, most of which assume one charger
per vehicle or give no depot-level guarantee for vehicles that must leave on time. Online methods use MPC
\cite{frendo2019realtime,tang2016model,lee2021adaptive}, flexibility-based priority ranking
\cite{lokhande2024realtime}, admission control \cite{alinia2020online}, stochastic MPC \cite{pozzi2022stochastic}
or approximate dynamic programming \cite{mahyari2025electric}.

Three gaps remain for rostered depots. (i) Priority is an exogenous flag, set without checking that the power
allocation can deliver every flagged vehicle's mission. (ii) Queue rules preempt only non-priority
vehicles; when every charger holds a priority vehicle there is no rule for which one may leave. (iii)
Receding-horizon controllers re-solve at every step, with no link from optimiser predictions to queue
decisions and no support for the dispatcher who must resolve overloads. AQPS addresses them with three
contributions:
\begin{enumerate}
  \item \textbf{Roster-derived priority with an optimiser feasibility test (Layer~0) and mission-safe
  charger allocation (Layer~1).} Vehicles are ranked by the Apparent Tardiness Cost (ATC) index
  \cite{vepsalainen1987priority,valente2007heuris}; the largest prefix for which the Layer-2 quadratic program
  (QP) meets every mission as a hard constraint becomes the priority set. Layer~1 preempts just in time and
  releases a priority vehicle only after its mission energy is delivered. Proposition~\ref{prop:guarantee}
  gives sufficient conditions for the mission guarantee.
  \item \textbf{Event-triggered MPC coupled to admission (Layer~2)} \cite{heemels2012introduction}: a predicted
  priority shortfall triggers Layer~0.
  \item \textbf{An operator impact preview} that reports the cost of alternative priority sets and executes
  unless vetoed \cite{parasuraman2000model,parasuraman2010complacency}.
\end{enumerate}
The evaluation uses real AEMO rooftop-PV data and seven methods, including a perfect-information
mixed-integer benchmark, and reports the trade-off with non-priority service.
Sections~\ref{sec:formulation}--\ref{sec:conclusion} give the formulation, AQPS, the case study, results and
conclusions.

% =============================================================================
\section{Problem formulation}
\label{sec:formulation}

\begin{figure}[!t]
  \centering
  \includegraphics[width=\columnwidth]{fig1_depot.pdf}
  \caption{Depot model: $M$ EVSEs behind one transformer with PV and BESS. Vehicles are connected ($\Cset$)
  or wait in the priority queue $Q^p_k$ or the non-priority queue $Q^n_k$. The diagram was drawn in TikZ with
  AI assistance (see the declaration).}
  \label{fig:depot}
\end{figure}

\subsection{Depot model}
A depot has $M$ identical Level-2 AC EVSEs, PV and a BESS behind one transformer (Fig.~\ref{fig:depot}). Time
steps have length $\Delta t$; $k$ is the current step. Vehicles present form $\mathcal{I}_k$, the connected
subset is $\Cset$ ($|\Cset|\le M$). Vehicle $i$ arrives at $a_i$, departs at $d_i$, has maximum charging
power $\bar{x}_i$ and arrival energy $E^{\mathrm{arr}}_i$. We assume that vehicles outnumber EVSEs; that the
roster gives each vehicle's next-trip energy $E^{\mathrm{trip}}_i$, return energy $E^{\mathrm{ret}}_i$,
further trips today, next-day departure and criticality weight $w_i$, and that disruptions change these
during the day; that a day-ahead PV forecast $P^{\mathrm{PV}}_t$ and a BESS schedule $P^{\mathrm{BESS}}_t$
(positive when discharging) are known while realised PV differs; and that Level-2 charging is in the
constant-current region, so degradation is not modelled. Each vehicle has a \emph{mission energy} $\Mis_i$
(Section~\ref{sec:qp}), needed for its next trip and return with reserve $r_i$, and a \emph{session energy}
$\Ses_i\ge\Mis_i$ the operator would like it to receive; in the QP both are remaining quantities. Priority
vehicles form the set $\Sset$ chosen by Layer~0.

\subsection{Layer-2 optimisation problem}
\label{sec:qp}
At epoch $k$ the powers $x_{i,t}$ of the connected vehicles over $\mathcal{T}_k=\{k,\dots,k+H_k-1\}$, with
$H_k=\max_{i\in\Cset}(d_i-k)$ reaching the latest departure, solve
\begin{equation}
  \min_{x,\,g}\ \ \alpha_1 H^{\mathrm{EC}}+\alpha_2 H^{\mathrm{NC}}+\alpha_3 H^{\mathrm{QC}},
  \label{eq:obj}
\end{equation}
with energy cost of net import, aggregate non-completion and a quadratic regulariser:
\begin{align}
  H^{\mathrm{EC}}&=\sum_{t\in\mathcal{T}_k} C_t\,g_t\,\Delta t, \notag\\
  &\qquad g_t\ge 0,\ \ g_t\ge P_t-P^{\mathrm{PV}}_t-P^{\mathrm{BESS}}_t, \label{eq:hec}\\
  H^{\mathrm{NC}}&=\Bigl(\sum_{i\in\Cset}\bigl(\Ses_i-\textstyle\sum_{t}x_{i,t}\Delta t\bigr)\Bigr)^{2}, \label{eq:hnc}\\
  H^{\mathrm{QC}}&=\sum_{t\in\mathcal{T}_k}P_t^{2},\qquad P_t=\textstyle\sum_{i\in\Cset}x_{i,t}. \label{eq:hqc}
\end{align}
$C_t$ is the ToU price; at the optimum $g_t=\max(0,P_t-P^{\mathrm{PV}}_t-P^{\mathrm{BESS}}_t)$, so only
imported energy is charged. For all $t\in\mathcal{T}_k$:
\begin{align}
  &|\mathbb{K}_l|\,I_t\le c_l,\quad l=1,\dots,L, \label{eq:branch}\\
  &P_t\le\Psite, \label{eq:site}\\
  &0\le x_{i,t}\le\bar{x}_i\ \text{for } a_i\le t<d_i,\ \ x_{i,t}=0\ \text{otherwise}, \label{eq:rate}\\
  &x_{i,t}\ge x^{\min}_{i,t},\quad i\in\Cset\cap\Sset, \label{eq:floor}\\
  &\textstyle\sum_{t}x_{i,t}\Delta t\le\Ses_i,\quad i\in\Cset, \label{eq:cap}\\
  &\textstyle\sum_{t}x_{i,t}\Delta t\ge\Mis_i,\quad i\in\Cset\cap\Sset, \label{eq:mission}
\end{align}
where
\begin{equation}
  \Mis_i=\min\!\bigl(\Ses_i,\ \max(0,\ E^{\mathrm{trip}}_i+E^{\mathrm{ret}}_i+r_i-E^{\mathrm{arr}}_i)\bigr).
  \label{eq:emission}
\end{equation}
Eq.~\eqref{eq:branch} is the phase-aware feeder limit of \cite{lee2021adaptive}; with the single feeder
studied here ($L=1$) it reduces to the gross site limit Eq.~\eqref{eq:site}, so a PV shortfall raises cost but
never causes infeasibility. Eq.~\eqref{eq:floor} gives a connected priority vehicle a front-loaded minimum
rate until its mission is covered: with $\xpri_i=\min(3.84\,\mathrm{kW},\bar{x}_i)$ (16\,A at 240\,V) and
$n_i=\lfloor\Mis_i/(\xpri_i\Delta t)\rfloor$, $x^{\min}_{i,t}=\xpri_i$ for $n_i$ steps, then the remainder
$(\Mis_i-n_i\xpri_i\Delta t)/\Delta t$ for one step, then zero. Eq.~\eqref{eq:cap} caps every vehicle at its
session energy and keeps the bracket in Eq.~\eqref{eq:hnc} non-negative; Eq.~\eqref{eq:mission} is hard for
the priority set, which Layer~0 admits only when it is satisfiable.

\textbf{Weights.} $(\alpha_1,\alpha_2,\alpha_3)=(1,100,10^{-5})$ follow from marginal values. Moving 1\,kW of
import between steps changes $\alpha_1H^{\mathrm{EC}}$ by $(C_t-C_{t'})\Delta t\ge(0.21-0.18)\times0.25=0.0075$\,\$
across tiers and $\alpha_3H^{\mathrm{QC}}$ by about $2\alpha_3(P_t-P_{t'})$, so the regulariser cannot reverse a
tariff preference while $|P_t-P_{t'}|<375$\,kW, which holds for every site limit of the main study
($\le300$\,kW); within a tier it spreads load. The marginal value of aggregate shortfall, $2\alpha_2S$ per kWh,
exceeds the highest tariff (0.32\,\$/kWh) once $S>1.6$\,Wh: sessions are completed whenever feasible and the
tariff decides \emph{when}. Section~\ref{sec:res_sens} confirms both.

% =============================================================================
\section{AQPS}
\label{sec:aqps}

\begin{figure*}[!t]
  \centering
  \includegraphics[width=\textwidth]{fig2_architecture.pdf}
  \caption{AQPS closed loop. The event detector selects the layers that run (L0--L2); without an event the
  cached Layer-2 plan is applied. $^\ast$Shortfall: the plan falls short of a connected priority vehicle's
  mission, or a waiting priority vehicle's slack reaches $\delta$; it triggers Layer~0 (dashed). Escalations
  and injections go to the impact preview, whose choice sets the priority set. Drawn in TikZ with AI assistance
  (see the declaration).}
  \label{fig:arch}
\end{figure*}

Layer~0 decides which vehicles are priority, Layer~1 which are connected and Layer~2 how much power each
receives (Fig.~\ref{fig:arch}).

\subsection{Layer 0: priority set}
\label{sec:layer0}
\textbf{Candidates} are present vehicles rostered as priority or with $w_i\ge w_{\min}=2$ that still need
mission energy and were not safely released. Connected priority candidates below their mission are
\emph{committed} and never demoted. \textbf{Urgency} is the ATC index \cite{vepsalainen1987priority}
\begin{equation}
  \pi_i(t)=\frac{w_i}{p_i(t)}\exp\!\left(-\frac{\max(0,s_i(t))}{k_{\mathrm{ATC}}\,\bar{p}(t)}\right),
  \label{eq:atc}
\end{equation}
with processing time $p_i=\max(\Delta t,\Mis_i/\bar{x}_i)$, slack $s_i=(d_i-t)\Delta t-p_i$, mean $\bar{p}$ over
present vehicles and $k_{\mathrm{ATC}}=1.5$ \cite{valente2007heuris}. \textbf{Feasibility test.} A candidate
failing the dwell test $\Mis_i\le\bar{x}_i(d_i-t)\Delta t$ is escalated at once, so that it cannot block the
ranking. A prefix of the remaining ranking is feasible if the Layer-2 problem with hard
Eq.~\eqref{eq:mission} over the committed vehicles plus the first $M-|\text{committed}|$ new candidates is
feasible with no predicted shortfall. The full prefix is tested first and bisection is used only if it
fails (at most $\lceil\log_2m\rceil+1$ solves for $m$ candidates). The test is exact for vehicles that can be
connected now and uses the dwell test, a necessary condition, for queued ones. Best-subset selection is
knapsack-like and NP-hard in general \cite{kellerer2004knapsack}; the perfect-information benchmark bounds
the loss of the greedy prefix. Excluded candidates are escalated.

\subsection{Layer 1: charger allocation with safe release}
\label{sec:layer1}
Layer~1 keeps the active priority queue $Q^{\mathrm{act}}_k$, the priority waiting queue $Q^p_k$ (ATC order)
and the non-priority queue $Q^n_k$ (arrival order); free EVSEs go to $Q^p_k$, then $Q^n_k$. When all $M$ are
busy, a waiting priority vehicle displaces nobody while its mission slack
$\sigma_i=(d_i-k)\Delta t-\Mis_i/\bar{x}_i$ exceeds $\delta=2$\,h (just-in-time preemption). Once
$\sigma_i\le\delta$:
\begin{enumerate}
  \item \emph{Tier~1:} disconnect the connected non-priority vehicle with the largest laxity
  $(d_j-k)\Delta t-\Ses_j/\bar{x}_j$.
  \item \emph{Tier~2, safe release:} disconnect a connected priority vehicle $j$ satisfying, in kWh,
  \begin{align}
    &E^{\mathrm{cur}}_j\ge E^{\mathrm{trip}}_j+E^{\mathrm{ret}}_j+r_j\ \text{and}\ \Mis_j\ \text{delivered}, \label{eq:sr_a}\\
    &E^{\mathrm{fut}}_j\le E^{\mathrm{cur}}_j-(E^{\mathrm{trip}}_j+E^{\mathrm{ret}}_j+r_j), \label{eq:sr_b}\\
    &\bar{x}_j(d^{\mathrm{next}}_j-a^{\mathrm{ret}}_j)\Delta t\ge E^{\mathrm{next}}_j-\varepsilon_j, \label{eq:sr_c}
  \end{align}
  with battery energy $E^{\mathrm{cur}}_j$, further trip energy today $E^{\mathrm{fut}}_j$, surplus
  $\varepsilon_j$ (right side of Eq.~\eqref{eq:sr_b}) and Eq.~\eqref{eq:sr_c} checking that the overnight window
  from evening return $a^{\mathrm{ret}}_j$ to next-day departure $d^{\mathrm{next}}_j$ can deliver the next
  mission $E^{\mathrm{next}}_j$ (skipped for vehicles staying overnight). The eligible vehicle with the largest
  surplus is released and becomes non-priority.
  \item \emph{Tier~3, guarded rotation:} rotate out the first vehicle $j$ of $Q^{\mathrm{act}}_k$ with
  $\pi_{\mathrm{new}}>1.02\,\pi_j$ that has reached its mission level or has slack above one step; it returns
  to $Q^p_k$ if its mission is unmet. In operation this is the dispatcher's decision through the preview; in
  simulation Tier~3 stands in for it, lies outside Proposition~\ref{prop:guarantee} and is counted separately.
  \item \emph{Tier~4:} disconnect nobody and open one escalation episode.
\end{enumerate}
Section~\ref{sec:res_sens} shows why $\delta$ matters.

\subsection{Layer 2: event-triggered MPC and operator preview}
\label{sec:et}
Layer~2 solves Eqs.~\eqref{eq:obj}--\eqref{eq:emission} over $\Cset$ and caches the plan $X^\ast$. Without an
event the next column is applied (clipped to the remaining session energy and $\bar{x}_i$); with an event
the layers of Table~\ref{tab:events} run first. The \emph{shortfall} event re-opens admission when optimiser
predictions show a priority mission at risk. If Eq.~\eqref{eq:mission} is infeasible, Layers~0--1 run once
more; if it is still infeasible, Eq.~\eqref{eq:mission} becomes a weighted penalty and the event is logged.
AQPS-RT re-solves every step, AQPS-RT-ET only on events.

\begin{table}[!htbp]
\centering\footnotesize
\caption{Event routing in AQPS-RT-ET ($\checkmark$: layer runs).}
\label{tab:events}
\begin{tabular}{lccc}
\toprule
Event & Layer 0 & Layer 1 & Layer 2 \\
\midrule
Arrival & $\checkmark$ & $\checkmark$ & $\checkmark$ \\
Operator injection (+ preview) & $\checkmark$ & $\checkmark$ & $\checkmark$ \\
Predicted shortfall / slack $\le\delta$ & $\checkmark$ & $\checkmark$ & $\checkmark$ \\
Departure or session complete & & $\checkmark$ & $\checkmark$ \\
PV deviation $>\tau_{\mathrm{pv}}$ (3.5\,kW) & & & $\checkmark$ \\
Periodic, $T_{\max}$ (60\,min) & & & $\checkmark$ \\
\bottomrule
\end{tabular}
\end{table}

On an operator injection AQPS solves three what-if problems: the recommended set, \emph{keep all urgent}
(adds the most urgent non-priority vehicle) and \emph{downgrade} (removes the least urgent priority vehicle,
preferably one whose mission is met), each reporting horizon cost and departure energies. The recommended
option executes unless vetoed \cite{parasuraman2000model}, so charging never waits; previews are limited to
injections to avoid alert fatigue \cite{parasuraman2010complacency}. In simulation the dispatcher's choice is
scripted. Algorithm~\ref{alg:aqps} summarises one step.

\begin{algorithm}[!t]
\caption{One AQPS-RT-ET control step $k$}
\label{alg:aqps}
\footnotesize
\begin{algorithmic}[1]
\State detect events (Table~\ref{tab:events}); \textbf{if} none: apply column $k-k_0$ of $X^\ast$, \Return
\If{arrival, injection or shortfall}
  \State Layer 0: escalate dwell-infeasible candidates; rank by Eq.~\eqref{eq:atc}; test full prefix, bisect if needed
\EndIf
\If{arrival, injection, shortfall or departure}
  \State Layer 1: admit from $Q^p_k$; if full and $\sigma_i\le\delta$: Tiers 1--3, else Tier 4; fill from $Q^n_k$
\EndIf
\If{injection} \State impact preview (3 what-if solves); apply scripted or recommended option \EndIf
\State Layer 2: solve; if infeasible re-run Layers 0--1, then logged soft solve; store $X^\ast$, $k_0\leftarrow k$
\end{algorithmic}
\end{algorithm}

\subsection{Properties}
\label{sec:theory}

\begin{proposition}[Convexity and unique aggregate profile]
\label{prop:convex}
For given $\Cset$ and $\Sset$, problem \eqref{eq:obj}--\eqref{eq:emission} is a convex QP in $(x,g)$. If it is
feasible (always so when $\Cset\cap\Sset=\emptyset$), an optimum exists, the optimal value is unique and all
optima share the aggregate profile $P^\ast_t$ and net import $g^\ast_t$; the split among vehicles need not be
unique.
\end{proposition}
\begin{proof}[Proof sketch]
The constraints are affine and bound $x$; with $\Cset\cap\Sset=\emptyset$, $x=0$ is feasible. As
$\alpha_1C_t>0$, optimal $g_t=\max(0,P_t-P^{\mathrm{PV}}_t-P^{\mathrm{BESS}}_t)$, and the objective becomes
$\phi(P)=\alpha_1\sum_tC_t\max(0,P_t-P^{\mathrm{PV}}_t-P^{\mathrm{BESS}}_t)\Delta t+\alpha_2(\sum_i\Ses_i-\Delta t\sum_tP_t)^2
+\alpha_3\sum_tP_t^2$, continuous and strictly convex in $P$ ($\alpha_3>0$), so a minimum exists on the
compact feasible set. Two optima with $P^{(1)}\ne P^{(2)}$ would have a feasible midpoint with lower $\phi$;
hence $P^\ast$ and $g^\ast$ are unique. Exchanging $\pm\epsilon$ between two vehicles at two steps keeps all
row and column sums, so $x^\ast$ need not be unique.
\end{proof}
A Layer-2 problem has $(|\Cset|+1)H_k\le4\,320$ variables for $|\Cset|\le35$, $H_k\le120$, and is solved in
polynomial time by interior-point methods \cite{boyd2004convex}.

\begin{proposition}[Mission guarantee for admitted vehicles]
\label{prop:guarantee}
If (A1) a committed vehicle's departure and mission do not change after commitment, (A2) no committed vehicle
is rotated out by Tier~3, and (A3) arithmetic is exact, then every vehicle committed by Layer~0 departs with
at least $\Mis_i$.
\end{proposition}
\begin{proof}[Proof sketch]
(i) Layer~0 commits vehicles only if hard Eq.~\eqref{eq:mission} is feasible for the committed set and the new
vehicles, and never demotes. (ii) Tier~1 disconnects only non-priority vehicles and Tier~2 only vehicles with
$\Mis_j$ delivered; with (A2) no committed vehicle leaves its charger early. (iii) If $X^\ast$ is feasible at
$k$, its tail is feasible at $k+1$: Eqs.~\eqref{eq:site}--\eqref{eq:cap} shift in time, the remaining mission
falls by the energy delivered, and the recomputed floor Eq.~\eqref{eq:floor} is pointwise no larger than the
shifted one; removing vehicles relaxes Eq.~\eqref{eq:site} and new non-priority vehicles are feasible with
$x=0$. (iv) Applied power equals the plan unless the session is complete, and $\Ses_i\ge\Mis_i$. By induction
Eq.~\eqref{eq:mission} holds up to departure.
\end{proof}
The guarantee covers the automatic tiers and admitted vehicles. Admission requires (C1) the mission to fit the
dwell at full power, (C2) a free or freeable (Tier~1--2) charger when $\sigma_i\le\delta$, and (C3) all
committed missions to fit under $\Psite$; otherwise the vehicle is escalated (Section~\ref{sec:res_overload}).
Disruptions violating (A1) are handled by re-running Layer~0 and, if needed, the logged soft solve; every run
reports these counts, so (A3) can be checked.

% =============================================================================
\section{Case study}
\label{sec:setup}

\subsection{Depot, fleet and scenarios}
Two depots are studied: 45 vehicles with 16 EVSEs and $\Psite=150$\,kW, and 90 vehicles with 35 EVSEs and
300\,kW; a 180-vehicle depot (60 EVSEs, 600\,kW) tests scalability. EVSEs are 22\,kW Level-2 AC. Vehicles are
light (45\,kWh, 11\,kW) and medium (68\,kWh, 22\,kW) vans \cite{vales2026emergent,fiori2018modelling} with
synthetic rosters: arrival energy 12--25\,\% of capacity, $r_i=0.1\,C_i$, next trip sized so that $\Mis_i$ is
70--85\,\% of $\Ses_i$, next-day departure 05:00--09:00 and $w_i$ of 2.5--4.5 (priority) or 0.5--1.4. A run
covers 06:00 on day $D$ to 12:00 on $D+1$ (120 steps of 15\,min, AEST). The site has 20\,kW of PV and a
10\,kWh BESS charging at 2.5\,kW over 10:00--14:00 and discharging about 1.5\,kW over 15:00--21:00 (90\,\%
round trip); BESS charging counts as load. PV and BESS are deliberately small, as on a retrofitted depot roof
\cite{ratnam2015scheduling,teng2020technical}, so that results reflect charger and transformer contention.
The ToU tariff \cite{nimalsiri2021coordinated} is 0.18\,\$/kWh off-peak (22:00--07:00), 0.21 shoulder
(07:00--15:00, 21:00--22:00) and 0.32 peak (15:00--21:00).

Scenarios (Table~\ref{tab:scenarios}) increase contention, defined before any run. S1: 27\,\% priority,
arrivals spread over the day. S2: an early cohort of high-demand non-priority vans holds the chargers when
priority vans arrive at 10:00--13:00. S3: the blocking cohort arrives at 13:00--15:00 and priority vans arrive
at 15:00--18:00, inside the peak.

\input{tab_scenarios}

\subsection{PV data}
PV is AEMO NEMWeb data for VIC1 \cite{aemo2025mms}: \texttt{ROOFTOP\_PV\_ACTUAL} (type MEASUREMENT) and
\texttt{ROOFTOP\_PV\_FORECAST}, using for day $D$ the latest run issued by 12:00 on $D-1$. 30-min values are
held for both 15-min steps and scaled to 20\,kW by the maximum VIC1 actual over July 2024--June 2025
(\RPVNormMW\,MW). The \RCandidateDays{} complete weekdays were ranked by daily normalised mean absolute error
(nMAE, 06:00--20:00, \% of 20\,kW) and the days at the 10th, 25th, 50th, 75th, 90th and 97.5th percentiles
selected (nMAE \RNmaeMin--\RNmaeMax\,\%). The median day, \RPVRefDate{} (nMAE \RPVRefNmae\,\%), drives all
experiments except E4. Optimising methods receive the forecast with a persistence correction decaying over
2\,h.

\subsection{Methods and metrics}
\label{sec:methods}
All methods share chargers, site limit, data and simulator.
\textbf{LLF} and \textbf{EDF} \cite{lee2021adaptive} give EVSEs and power by laxity or deadline; a connected
vehicle keeps its EVSE until it finishes or departs. \textbf{MPC-only} solves the Layer-2 problem every step
with first-come first-served admission and no priority set. \textbf{Offline-Opt.} solves the same problem over
the window with perfect knowledge (without Eq.~\eqref{eq:floor}), the EVSE count enforced by binary variables
($\sum_iz_{i,t}\le M$, $x_{i,t}\le\bar{x}_iz_{i,t}$) and Eq.~\eqref{eq:mission} hard for rostered priority
vehicles, as a mixed-integer QP (1\,\% gap target, gap bounded by the relaxation). \textbf{AQPS-static} keeps
day-ahead flags, Tier~1 only and Layer~2 every step. \textbf{AQPS-RT} and \textbf{AQPS-RT-ET} are the full
method.

\emph{Mission sufficiency} is the share of rostered priority vehicles departing with at least $\Mis_i$
(0.1\,kWh tolerance); \emph{non-priority fulfilment} is delivered over requested non-priority session energy,
at the end of the window and at 22:00; \emph{unit cost} is grid cost per kWh delivered; \emph{solves} counts
all optimiser solves. The simulator uses MOSEK \RMosekVersion{} \cite{mosek2025} via CVXPY \RCvxpyVersion{}
\cite{diamond2016cvxpy} without solver fallback; all results come from commit \texttt{\RSimCommit} on an
\RCPUCores-core Intel processor, and every number here is generated from the committed results
\cite{uddin2026aqps}.

% =============================================================================
\section{Results}
\label{sec:results}

\subsection{Main comparison (E1)}
\label{sec:res_main}

\input{tab_e1_master}

\begin{figure*}[!t]
  \centering
  \includegraphics[width=\textwidth]{fig3_service.pdf}
  \caption{Mission sufficiency (top) and non-priority energy fulfilment (bottom), E1. Offline-Opt.\ (hatched)
  has perfect information.}
  \label{fig:service}
\end{figure*}

In S1 (Table~\ref{tab:e1}, Fig.~\ref{fig:service}) every method meets every mission; the optimising methods
shift non-priority charging off-peak, lowering unit cost by \RUnitSaveMinSOne--\RUnitSaveMaxSOne\,\% against LLF
(AQPS-RT-ET \RUnitCostRTETSOneFortyFive{} and \RUnitCostRTETSOneNinety\,\$/kWh vs \RUnitCostLLFSOneFortyFive{}
and \RUnitCostLLFSOneNinety), at the price of later delivery: at 22:00 AQPS-RT-ET has delivered
\RNonPriTwoTwoRTETSOneFortyFive{} and \RNonPriTwoTwoRTETSOneNinety\,\% of non-priority energy, LLF
\RNonPriTwoTwoLLFSOneFortyFive{} and \RNonPriTwoTwoLLFSOneNinety\,\%. In S2 LLF still meets every mission,
EDF misses one at 90 vehicles and MPC-only meets \RMissionMPCSTwoFortyFive{} and \RMissionMPCSTwoNinety\,\%;
AQPS meets all, with less non-priority energy than LLF (\RNonPriRTETSTwoFortyFive{} and
\RNonPriRTETSTwoNinety\,\% vs \RNonPriLLFSTwoFortyFive{} and \RNonPriLLFSTwoNinety\,\%) at a
\RUnitSaveMinSTwo--\RUnitSaveMaxSTwo\,\% lower unit cost, so S2 trades non-priority energy for cost.

S3 is the case AQPS targets. With priority vans arriving in the peak while the blocking cohort holds the
chargers, LLF meets \RMissionLLFSThreeFortyFive{} and \RMissionLLFSThreeNinety\,\% of missions, EDF
\RMissionEDFSThreeFortyFive{} and \RMissionEDFSThreeNinety\,\% and MPC-only \RMissionMPCSThreeFortyFive{} and
\RMissionMPCSThreeNinety\,\%; all AQPS variants meet every mission at a unit cost
\RUnitSaveMinSThree--\RUnitSaveMaxSThree\,\% below LLF. Non-priority fulfilment falls to
\RNonPriRTETSThreeFortyFive{} and \RNonPriRTETSThreeNinety\,\% (LLF \RNonPriLLFSThreeFortyFive{} and
\RNonPriLLFSThreeNinety\,\%), but the perfect-information benchmark, which must also meet every mission,
reaches only \RNonPriOfflineSThreeFortyFive{} and \RNonPriOfflineSThreeNinety\,\%: most of the loss is the
price of the missions, and AQPS is within \RSThreeOfflineGapMin--\RSThreeOfflineGapMax{} points of the bound.
The benchmark solved to a gap of at most \ROfflineGapMax\,\% within \ROfflineSolveMaxS\,s. Without disruptions
the day-ahead flags are correct, so AQPS-static also meets every mission; its value is tested in E2, E3 and
E7.

\begin{figure*}[!t]
  \centering
  \includegraphics[width=\textwidth]{fig4_cost_tiers.pdf}
  \caption{Grid energy cost by tariff tier, E1, S2 and S3; numbers above bars are unit costs (\$/kWh).}
  \label{fig:cost}
\end{figure*}

Total cost (Fig.~\ref{fig:cost}) differs mostly because methods deliver different energy; in S3 most priority
energy falls in the peak for every method. Fig.~\ref{fig:gantt} shows the mechanism: under LLF the priority
vans wait while earlier non-priority vans keep their chargers and leave \RGanttShortMin--\RGanttShortMax\,kWh
short; under AQPS-RT-ET they preempt the non-priority vans with most slack when their own slack reaches
$\delta$. The preempted vans leave with \RGanttNPAqpsMin--\RGanttNPAqpsMax\,\% of their session energy
(\RGanttNPLLFMin--\RGanttNPLLFMax\,\% under LLF). The aggregate penalty Eq.~\eqref{eq:hnc} is indifferent to
which vehicle receives energy, so under a binding site limit the QP spreads power, while LLF gives full power
and frees chargers sooner; a preemption then ends a session planned to finish later.

\begin{figure*}[!t]
  \centering
  \includegraphics[width=\textwidth]{fig5_gantt.pdf}
  \caption{Charging of six vehicles, 90 vehicles, S3: (a) AQPS-RT-ET, (b) LLF; the three priority vehicles (P)
  with the largest LLF shortfall and the three non-priority vehicles (NP) most often preempted by
  AQPS-RT-ET.}
  \label{fig:gantt}
\end{figure*}

\subsection{Scripted disruptions (E2)}
\label{sec:res_e2}

\input{tab_e2_cases}

E2 injects four disruptions into the 90-vehicle S2 depot (Table~\ref{tab:e2}): (1) at 12:00 a priority van
breaks down and a standby van is promoted with a 32\,kWh mission and a 15:30 departure; (2) at 13:00 a
driver absence delays a priority van by 2\,h; (3) at 11:30 a route moves from a priority to a non-priority
van and the dispatcher keeps all urgent vans; (4) at 12:00 realised PV drops to 25\,\% of forecast for 2\,h.
AQPS-static misses the promoted van's mission in case~1 (\RCaseOneMissionStatic\,\%); AQPS-RT-ET meets every
mission in every case, while in case~3 LLF misses one mission and EDF three. The PV drop changes only cost.

\begin{figure*}[!t]
  \centering
  \includegraphics[width=\textwidth]{fig6_swimlane.pdf}
  \caption{Case 1, AQPS-RT-ET: charger state of the promoted van (EV~009$^\ast$) and of vehicles disconnected
  within 4\,h of the injection.}
  \label{fig:swimlane}
\end{figure*}

In Fig.~\ref{fig:swimlane} EV~009 charges as non-priority until the injection, is promoted and reaches its
mission shortly before departure. When waiting priority vans run out of slack, vans that have reached their
missions are safely released (squares just after the mission markers) and one van is rotated and
readmitted. The preview at the injection (Table~\ref{tab:preview}) shows that keeping the next most urgent van
costs nothing extra, while a downgrade would cost \$\RCaseOneDowngradeDelta{} more.

\input{tab_e2_preview}

\subsection{Random disruptions (E3)}
\label{sec:res_e3}

\input{tab_e3_montecarlo}

E3 perturbs the 90-vehicle S3 depot over \REThreeSeeds{} seeds: arrivals and departures shift by up to
$\pm1$\,h, trip and session energies scale by $[0.88,1.15]$, and each daytime non-priority van is promoted
during its stay with probability 0.08 (Table~\ref{tab:e3}). AQPS-RT-ET meets
\REThreeMissionRTET$\,\pm\,$\REThreeMissionCIRTET\,\% of missions; its misses, in \REThreeMissSeedsRTET{}
seeds, are each a van promoted when its mission could no longer be delivered at full power. AQPS-static,
which ignores promotions, meets \REThreeMissionStatic\,\%, MPC-only \REThreeMissionMPC, EDF
\REThreeMissionEDF{} and LLF \REThreeMissionLLF\,\%. Unit costs differ by at most \REThreeUnitSpreadPct\,\%;
non-priority fulfilment is \REThreeNonPriRTET\,\% for AQPS-RT-ET and \REThreeNonPriLLF\,\% for LLF.

\subsection{PV forecast error (E4)}
\label{sec:res_e4}

\begin{figure*}[!t]
  \centering
  \includegraphics[width=\textwidth]{fig7_aemo.pdf}
  \caption{(a)--(c) AEMO VIC1 rooftop PV scaled to 20\,kW, actual and day-ahead forecast, on the 10th, 50th
  and 97.5th percentile days of forecast error. (d)--(e) Unit cost against daily nMAE, S2.}
  \label{fig:aemo}
\end{figure*}

E4 repeats S2 on the six AEMO days (Fig.~\ref{fig:aemo}). With a gross site limit PV affects cost, not
feasibility, so every method's service metrics are identical on all days: E4 is a cost-sensitivity study.
Unit cost follows the day's PV energy more than its forecast error (on the worst day the forecast exceeded
realised PV energy by \RPVOverForecastPct\,\%), the ranking of methods never changes, and AQPS-RT-ET is
\REFourSaveMinFortyFive--\REFourSaveMaxFortyFive\,\% (45 vehicles) and
\REFourSaveMinNinety--\REFourSaveMaxNinety\,\% (90) cheaper per kWh than LLF. Regional AEMO data understate the
variability a single 20\,kW array would see.

\subsection{Event triggering and scalability (E5)}
\label{sec:res_et}

\input{tab_e5_et}

Against every-step AQPS-RT (Table~\ref{tab:e5a}), event triggering removes \RLTwoRedMin--\RLTwoRedMax\,\% of
Layer-2 solves and \RTotRedMin--\RTotRedMax\,\% of all solves; mission sufficiency is identical, non-priority
fulfilment differs by at most \RETMaxDeltaNonPri{} points and unit cost by at most \RETMaxDeltaCostMWh\,\$/MWh.
E1 Layer-2 solves take \REOneMeanMsMinRTET--\REOneMeanMsMaxRTET\,ms on average (95th percentile at most
\REOnePNinetyFiveMaxRTET\,ms) against a 15-min step. Raising $T_{\max}$ from 15 to 120\,min (90 vehicles, S2)
cuts Layer-2 solves from \RTmaxLTwoFirst{} to \RTmaxLTwoLast{} with unchanged outcomes, because events, not
the timer, drive re-planning.

\input{tab_e5_scal}

In S2 with 45--180 vehicles and 2--4 vehicles per charger (Table~\ref{tab:e5c}), LLF and AQPS-RT-ET meet every
mission. AQPS-RT-ET's non-priority fulfilment falls from \RScalNonPriRTETLowMin--\RScalNonPriRTETLowMax\,\% at
2:1--3:1 to \RScalNonPriRTETFortyFiveRatioFour--\RScalNonPriRTETOneEightyRatioFour\,\% at 4:1, while LLF keeps
\RScalNonPriLLFFortyFiveRatioFour--\RScalNonPriLLFOneEightyRatioFour\,\%; its unit cost is
\RScalSaveMin--\RScalSaveMax\,\% lower. In S2 priority vans arrive after the blocking cohort is served, so
AQPS's non-priority loss buys no service gain. At 180 vehicles mean solve time is
\RScalMeanMsMin--\RScalMeanMsMax\,ms (95th percentile $\le$\RScalPNinetyFiveMax\,ms).

\subsection{Sensitivity (E6)}
\label{sec:res_sens}

\begin{figure*}[!t]
  \centering
  \includegraphics[width=\textwidth]{fig8_sensitivity.pdf}
  \caption{Sensitivity of AQPS-RT-ET, 90 vehicles, S2: (a)--(b) $\alpha_3$, (c)--(d) reserve $r$; dashed
  lines mark defaults.}
  \label{fig:sens}
\end{figure*}

One parameter at a time on the 90-vehicle S2 depot (Fig.~\ref{fig:sens}): unit cost is flat for
$\alpha_3\le10^{-4}$ and rises from $10^{-3}$, where the regulariser outweighs the tariff
(Section~\ref{sec:qp}); $\alpha_2$ (10--1000) and $\tau_{\mathrm{pv}}$ (1--8\,kW) change nothing measurable. The
preemption threshold matters most: non-priority fulfilment is \RDeltaTwoNonPri\,\% for $\delta$ of 1--2\,h,
\RDeltaThreeNonPri\,\% at 3\,h and \RDeltaEagerNonPri\,\% with eager preemption (the earlier design), with no
mission gain. $k_{\mathrm{ATC}}$ (0.5--3, at four vehicles per charger) changes non-priority fulfilment by under
one point; at $k_{\mathrm{ATC}}=3$ one mission is missed. Reserves of 5--15\,\% of capacity keep every mission;
at 20\,\% \RReserveTwentyMission\,\% are met, with one committed-infeasible event and soft solve, the only one
in all AQPS-RT and AQPS-RT-ET runs.

\subsection{Priority overload (E7)}
\label{sec:res_overload}

\input{tab_e7}

E7 makes 30--90\,\% of the 45-vehicle S3 fleet priority, within the contended cohort (Table~\ref{tab:e7}). Up
to 50\,\% both AQPS variants meet every mission (LLF and EDF
\RLowShareMissionMinLLFEDF--\RLowShareMissionMaxLLFEDF\,\%, MPC-only
\RLowShareMissionMinMPC--\RLowShareMissionMaxMPC\,\%). Beyond that, demand exceeds 16 chargers: AQPS-static,
unable to change flags, finds Eq.~\eqref{eq:mission} infeasible
\ROverloadSoftMinStatic--\ROverloadSoftMaxStatic{} times and meets
\ROverloadMissionMinStatic--\ROverloadMissionMaxStatic\,\%, while AQPS-RT-ET admits what it can serve,
escalates the rest, needs no soft solve and meets \ROverloadMissionMinRTET--\ROverloadMissionMaxRTET\,\%
(one vehicle missed per run), against \ROverloadMissionMinLLF--\ROverloadMissionMaxLLF\,\% for LLF and
\ROverloadMissionMinMPC--\ROverloadMissionMaxMPC\,\% for MPC-only.

\subsection{Discussion}
\label{sec:discussion}
AQPS improves mission sufficiency when priority demand arrives while chargers are held by less urgent
vehicles (S3, E3, E7), at an equal or lower cost per kWh; when arrival order already favours priority vehicles
(S1, S2, E5c), LLF meets the missions and serves non-priority vehicles better. The non-priority loss has an
unavoidable part, measured by the perfect-information benchmark, and an online part from preempting sessions
planned to finish later; a charger-turnover term in Eq.~\eqref{eq:hnc} could reduce the latter. In practice
AQPS sits between a fleet management system (rosters, weights, disruptions) and a charging station management
system (setpoints, e.g.\ OCPP~2.0.1 profiles). Limitations: the dispatcher is simulated by scripted choices and
Tier~3; rosters are synthetic; PV is AEMO's regional estimate scaled by the period maximum; each run covers one
30-h window; degradation, switching and communication delays are not modelled.

% =============================================================================
\section{Conclusion}
\label{sec:conclusion}
AQPS admits the largest priority set whose missions a convex QP can guarantee, preempts just in time,
releases priority vehicles only after their missions are met, re-plans on events and supports the dispatcher
with an impact preview. The QP has a unique optimal aggregate profile, and admitted priority vehicles depart
with their mission energy under stated assumptions. With real AEMO PV data AQPS met every mission in the
peak-stress scenario, where least-laxity-first met \RMissionLLFSThreeFortyFive{} and
\RMissionLLFSThreeNinety\,\%, at an equal or lower cost per kWh, kept
\ROverloadMissionMinRTET--\ROverloadMissionMaxRTET\,\% under priority overload and cut solves by
\RTotRedMin--\RTotRedMax\,\%. The price is lower non-priority energy, mostly inherent to meeting the missions.
Future work: charger-turnover awareness, joint depot--public charging, and validation with dispatchers and
hardware.

\section*{CRediT authorship contribution statement}
\textbf{Mohammad Kamar Uddin:} Conceptualization, Methodology, Software, Data curation, Formal analysis,
Writing -- original draft, Visualization. \textbf{Ali Moradi Amani:} Conceptualization, Methodology,
Writing -- review \& editing, Supervision. \textbf{Hui Song:} Methodology, Writing -- review \& editing,
Supervision.

\section*{Declaration of competing interest}
The authors declare that they have no known competing financial interests or personal relationships that
could have appeared to influence the work reported in this paper.

\section*{Funding}
% AUTHORS: confirm, or replace with: "This work was supported by [funder] [grant number]."
This research did not receive any specific grant from funding agencies in the public, commercial, or
not-for-profit sectors.

\section*{Data availability}
The simulator, configuration, AEMO download and preparation scripts, result summaries, and the scripts that
generate every figure, table and number in this paper are available at
\url{https://github.com/kamareee/aqps-ev-simulation} \cite{uddin2026aqps}.

\section*{Declaration of generative AI and AI-assisted technologies in the manuscript preparation process}
% AUTHORS: confirm that this statement describes your use accurately.
During the preparation of this work the authors used Claude (Anthropic) to review and write simulation and
post-processing code, to generate the figures and tables from the simulation results, to draw Figs.~1--2,
and to draft and edit the text. After using this tool, the authors reviewed and edited the content as needed
and take full responsibility for the content of the published article.

\bibliographystyle{elsarticle-num-names}
\bibliography{refs}

\end{document}
